\documentclass[12pt,a4paper, colorlinks = true,
            linkcolor = blue,
            urlcolor  = blue,
            citecolor = blue,
            anchorcolor = blue]{moderncv}

\moderncvtheme[blue]{classic}
\usepackage[utf8]{inputenc}  %Windows

%\usepackage[scale=0.975]{geometry}
\usepackage[top=2.5cm, bottom=2.5cm, left=2.5cm, right=2.5cm]{geometry}

% Set a very tall paper height (e.g. 200cm)
%\geometry{
%  paperwidth=21cm,
%  paperheight=200cm,
%  margin=2cm
%}
\usepackage{graphicx}

\firstname{Andrew}
\familyname{Winter}
\title{\small{ORCID: 0000-0002-7501-9801}\\\small{Curriculum Vitae -- \today } \vspace{-2ex}}
\address{Queen Mary University of London\\Mile End Road\\London E1 4NS, United Kingdom\\}
\email{\,andrew.winter@qmul.ac.uk}
\homepage{www.andrew-j-winter.com}
%\phone{}
%\fax{Votre Fax}
%\email{clement.schmitt@etud.univ-montp2.fr}
\extrainfo{}
%\photo[50pt]{photo.png}
%\quote{Objectif: Intégrer une équipe au sein d'un projet novateur en tant que développeur informatique.}
\makeatletter
\renewcommand*{\bibliographyitemlabel}{\@biblabel{\arabic{enumiv}}}
\makeatother

\usepackage{multibib}
\newcites{book,misc}{{Books},{Others}}

\nopagenumbers{}
\begin{document}
\maketitle
\vspace{-6ex}



\section{Overview}

I study how the environment surrounding forming stars shapes the protoplanetary discs and planets that emerge from them, connecting star-forming regions, disc evolution, stellar clusters and galactic-scale star formation to observed exoplanet demographics. My work is primarily theoretical, developed in close collaboration with observers and grounded in observational statistics.

\vspace{-4pt}
\section{Research Positions \& Funding Awards}

\cventry{2026--}{\textnormal{\textbf{Max-Planck-Gesellschaft (MPG) Partner Group Leader} -- MPIA partnership, Heidelberg, €100k}}{}{}{}{}
\cventry{05/25--}{\textnormal{\textbf{Royal Society University Research Fellowship}: \textit{`Environmental origins of exoplanet diversity' } --  Queen Mary University of London, £1.83M}}{}{}{}{}
\cventry{09/24--04/25}{\textnormal{Visiting researcher --   Max Planck Institute for Astronomy, Heidelberg}}{}{}{}{}
\cventry{10/23--04/25}{\textnormal{\textbf{Marie Skłodowska Curie Fellowship}: \textit{`Observational signatures of planet formation in irradiated discs' } --   Observatoire de la C\^{o}te d'Azur, €195k}}{}{}{}{}
\cventry{10/22--10/23}{\textnormal{PDRA -- Kinematic signatures of planet formation in protoplanetary discs, Observatoire de la C\^{o}te d'Azur}}{}{}{}{}
\cventry{6/22--7/22}{\textnormal{Visiting researcher -- Chemical inhomogeneities in globular clusters, Institute of Astronomy, University of Cambridge}}{}{}{}{}
\cventry{4/22}{\textnormal{\textbf{Paris Regional Fellowship} (declined), €130k}}{}{}{}{}
\cventry{3/20--3/22}{\textnormal{\textbf{Humboldt Fellowship} -- \textit{`Planet evolution in dense stellar clusters'}, Heidelberg University, €80k}}{}{}{}{}
\cventry{5/19--3/20}{\textnormal{PDRA -- \textit{`Photoevaporative disc winds around binary stars'}, University of Leicester}}{}{}{}{}
\cventry{9/18--11/18}{\textnormal{DFG funded visiting Ph.D. student  -- \textit{`Linking galactic-scale star formation physics to protoplanetary disc survival'}, Heidelberg University}}{}{}{}{}{}

\vspace{-12pt}
\section{Education}

\cventry{Sep 2026}{\textnormal{Certificate in Learning and Teaching (CILT), Queen Mary University of London}}{}{}{}{}

\cventry{2015--2019}{\textnormal{Ph.D., Institute of Astronomy, University of Cambridge \newline Graduation date: Oct 2019, Supervisor: Cathie Clarke \newline Thesis: The influence of stellar birth environment on protoplanetary disc evolution}}{}{}{}{}

\cventry{2011--2015}{\textnormal{MMathPhys (Master of Mathematics and Physics), University of Warwick \newline 1$^{\mathrm{st}}$ Class Honours, Supervisor: Andrew Levan \newline Dissertation: Host galaxies of core collapse supernovae}}{}{}{}{}

\vspace{-12pt}
\section{Publications}

I have published 74 peer-reviewed articles (18 as first author), with an h-index of 26 and 1847 total citations. A full publication list is available on \href{https://ui.adsabs.harvard.edu/user/libraries/PK0RWOWOTIKWfo-5Fck9sg}{NASA/ADS.}

\vspace{-4pt}
\section{Conferences \& Talks}

I have given 19 talks in departments across Europe, the US and China and at 42 international conferences since 2019. This includes 9 invited review talks at large international conferences.

\vspace{-4pt}
\section{Supervision \& Teaching}
\cventry{Supervision}{\textnormal{\textbf{PhD}: Oliver Brown (2025--)$^{*}$, Mika Kontiainen (2024--), Rossella Anania (2022--25), Daniele Fasano (2022--25) \newline \textbf{Masters}: Fahham Kurji (2026--)$^{*}$, James Wirth (2023--24)$^{*}$ \newline \textbf{Bachelors}: Linling Shuai (2022--23)$^{*}$ \newline \textbf{Summer}: Audrey Liu (2026)$^{*}$ \newline $^{*}$Main supervisor (others: co-supervision or informal)}}{}{}{}{}
\cventry{Teaching}{\textnormal{Physical Cosmology Lecturer/Module Co-organiser, QMUL (2025--present); Part II Supervisor -- General Relativity \& Astrophysical Fluid Dynamics, University of Cambridge (2015--19)}}{}{}{}{}

\vspace{-3pt}
\section{Community \& Collaborations}
I am an active member of several large collaborations: \href{https://www.exoalma.com/}{exoALMA} (ALMA large programme), \href{https://www.transitiondiscs.com/}{Transition Discs} (DFG research unit), \href{https://westerlund1survey.wordpress.com}{EWOCS} and \href{https://www2.mpia-hd.mpg.de/homes/ramirez/XUE_MCRT.html}{XUE}. I contribute to the community as a JWST Cycle 4 discussion panelist (previously Cycle 2/3 external panel) and an ALMA distributed peer reviewer, and as a grant reviewer for Fondazione Cariplo, the Humboldt Foundation and STFC, alongside regular refereeing for MNRAS, A\&A, AJ, ApJ and Nature Astronomy. I have organised several conferences and workshops, including the PLATO Theory group meeting (QMUL, 2025) and the exoALMA workshop (Nice, 2024).

\vspace{-3pt}
\section{Observing Time}
\cventry{NOEMA}{\textnormal{Principal Investigator, 3mm continuum/line survey -- 50.6h science time + overheads, Grade B (observations in progress)}}{}{}{}{}
\cventry{ALMA}{\textnormal{Co-Investigator on 9 accepted proposals (Cycles 2023--2026, Grades B/C) and Principal Investigator on 1 further programme, studying protoplanetary disc structure, chemistry and late-stage infall}}{}{}{}{}
\cventry{JWST}{\textnormal{Co-Investigator on 6 awarded programmes spanning Cycles 1--4 (PIs Bik, Kuhn, Guarcello $\times$2, Benisty and Ramirez-Tannus), studying externally irradiated and starburst environments}}{}{}{}{}

\vspace{-12pt}
\section{Outreach}
\cventry{2016--26}{\textnormal{Public talks and events, including a \href{https://www.rmg.co.uk/whats-on/online/think-space-lectures}{Think Space lecture} (2026) and Cambridge Science Festival talks and a gravity-teaching game, \href{https://github.com/ajw278/gravity_hero}{Gravity Hero} (2016--18)}}{}{}{}{}

\end{document}
